% GENERATED FILE. Edit canonical lessons, then run npm run generate:books.
% canonical-source-hash: fnv1a64:d6489c1d2f86a02f
% canonical-lessons: PA-C331-chacha, PA-C331-taia, PA-C331-mama, PA-C331-masi, PA-C331-masar

\chapter{Uncles and Aunts}
\label{ch:pa-a2-uncles-and-aunts}
\hlchaptermodality{331}

\begin{chapteropening}
\textbf{By the end of this chapter:} \emph{I can name my uncles and aunts on both sides.}

Complete the last lesson of chapter 331: I can name my uncles and aunts on both sides.
\end{chapteropening}

\section[\texorpdfstring{\pa{ਚਾਚਾ} (chāchā)}{ਚਾਚਾ (chāchā)}]{\pa{ਚਾਚਾ} (chāchā) — a father's younger brother}

\label{lesson:PA-C331-chacha}



\begin{quote}
Before the new one: say the Punjabi for a husband's younger brother's wife, then the Punjabi for a husband's elder brother's wife.
\end{quote}

\subsection*{You'll want to know: \pa{ਚਾਚਾ}}
\textbf{\pa{ਚਾਚਾ}} — \emph{chāchā} — \textquotedblleft{}a father's younger brother\textquotedblright{}.

\begin{quote}
\pa{ਮੇਰਾ} \pa{ਚਾਚਾ} \pa{ਖੇਤ} \pa{ਵਿੱਚ} \pa{ਕੰਮ} \pa{ਕਰਦਾ} \pa{ਹੈ।}
\end{quote}

\begin{quote}
\emph{merā chāchā khet vich kamm kardā hai.}
\end{quote}

\begin{quote}
My father's younger brother works in the fields.
\end{quote}

A masculine noun.

\begin{cousinweb}
\textbf{\pa{ਚਾਚਾ}} (\emph{chāchā}) is inherited from a reconstructed \emph{*cācca}, \textquotedblleft{}uncle\textquotedblright{}; CDIAL lists Punjabi \emph{cācā}, \textquotedblleft{}father's younger brother\textquotedblright{}, under it (CDIAL 4734).
\end{cousinweb}

\begin{grammarlens}[title={What you've built}]
One more word for this chapter.
\end{grammarlens}

\subsection*{Your turn}
\begin{itemize}
\raggedright
  \item \emph{Say it:} \emph{chāchā}
  \item \emph{Say it:} \emph{chāchā}, once more
  \item \emph{Say it:} \emph{merā chāchā khet vich kamm kardā hai.} — My father's younger brother works in the fields.
  \item \emph{Recall:} say the Punjabi for a husband's younger brother's wife, then the Punjabi for a husband's elder brother's wife, then say \emph{chāchā} again
\end{itemize}

\begin{tcolorbox}[breakable,colback=teal!4,colframe=teal!35!black,title={Before you move on}]
What does \pa{ਚਾਚਾ} mean? (\textquotedblleft{}A father's younger brother\textquotedblright{}.) Say it once more.
\end{tcolorbox}
\section[\texorpdfstring{\pa{ਤਾਇਆ} (tāiā)}{ਤਾਇਆ (tāiā)}]{\pa{ਤਾਇਆ} (tāiā) — a father's elder brother}

\label{lesson:PA-C331-taia}



\begin{quote}
Before the new one: say the Punjabi for a husband's elder brother's wife, then the Punjabi for a father's younger brother.
\end{quote}

\subsection*{You'll want to know: \pa{ਤਾਇਆ}}
\textbf{\pa{ਤਾਇਆ}} — \emph{tāiā} — \textquotedblleft{}a father's elder brother\textquotedblright{}.

\begin{quote}
\pa{ਮੇਰਾ} \pa{ਤਾਇਆ} \pa{ਸਾਡੇ} \pa{ਨਾਲ} \pa{ਰਹਿੰਦਾ} \pa{ਹੈ।}
\end{quote}

\begin{quote}
\emph{merā tāiā sāḍe nāl rahindā hai.}
\end{quote}

\begin{quote}
My father's elder brother lives with us.
\end{quote}

A masculine noun.

\begin{cousinweb}
\textbf{\pa{ਤਾਇਆ}} (\emph{tāiā}) is inherited from Sanskrit \emph{tātya}, \textquotedblleft{}fatherly\textquotedblright{}; CDIAL lists Punjabi \emph{tāiā}, \textquotedblleft{}father's elder brother\textquotedblright{}, under it (CDIAL 5759).
\end{cousinweb}

\begin{grammarlens}[title={What you've built}]
One more word for this chapter.
\end{grammarlens}

\subsection*{Your turn}
\begin{itemize}
\raggedright
  \item \emph{Say it:} \emph{tāiā}
  \item \emph{Say it:} \emph{tāiā}, once more
  \item \emph{Say it:} \emph{merā tāiā sāḍe nāl rahindā hai.} — My father's elder brother lives with us.
  \item \emph{Recall:} say the Punjabi for a husband's elder brother's wife, then the Punjabi for a father's younger brother, then say \emph{tāiā} again
  \item \emph{Write it:} copy \textbf{\pa{ਤਾਇਆ}} once, checking the whole word
\end{itemize}

\begin{tcolorbox}[breakable,colback=teal!4,colframe=teal!35!black,title={Before you move on}]
What does \pa{ਤਾਇਆ} mean? (\textquotedblleft{}A father's elder brother\textquotedblright{}.) Say it once more.
\end{tcolorbox}
\section[\texorpdfstring{\pa{ਮਾਮਾ} (māmā)}{ਮਾਮਾ (māmā)}]{\pa{ਮਾਮਾ} (māmā) — a mother's brother}

\label{lesson:PA-C331-mama}



\begin{quote}
Before the new one: say the Punjabi for a father's younger brother, then the Punjabi for a father's elder brother.
\end{quote}

\subsection*{You'll want to know: \pa{ਮਾਮਾ}}
\textbf{\pa{ਮਾਮਾ}} — \emph{māmā} — \textquotedblleft{}a mother's brother\textquotedblright{}.

\begin{quote}
\pa{ਗਰਮੀ} \pa{ਦੀ} \pa{ਛੁੱਟੀ} \pa{ਵਿੱਚ} \pa{ਅਸੀਂ} \pa{ਮਾਮੇ} \pa{ਦੇ} \pa{ਘਰ} \pa{ਗਏ।}
\end{quote}

\begin{quote}
\emph{garmī dī chhuṭṭī vich asīṁ māme de ghar gae.}
\end{quote}

\begin{quote}
In the summer holidays we went to my mother's brother's house.
\end{quote}

A masculine noun.

\begin{cousinweb}
\textbf{\pa{ਮਾਮਾ}} (\emph{māmā}) is inherited from Sanskrit \emph{māma}, \textquotedblleft{}uncle\textquotedblright{}, a nursery word; CDIAL lists Punjabi \emph{māmmā} under it (CDIAL 10055).
\end{cousinweb}

\begin{grammarlens}[title={What you've built}]
One more word for this chapter.
\end{grammarlens}

\subsection*{Your turn}
\begin{itemize}
\raggedright
  \item \emph{Say it:} \emph{māmā}
  \item \emph{Say it:} \emph{māmā}, once more
  \item \emph{Say it:} \emph{garmī dī chhuṭṭī vich asīṁ māme de ghar gae.} — In the summer holidays we went to my mother's brother's house.
  \item \emph{Recall:} say the Punjabi for a father's younger brother, then the Punjabi for a father's elder brother, then say \emph{māmā} again
  \item \emph{Write it:} copy \textbf{\pa{ਮਾਮਾ}} once, checking the whole word
\end{itemize}

\begin{tcolorbox}[breakable,colback=teal!4,colframe=teal!35!black,title={Before you move on}]
What does \pa{ਮਾਮਾ} mean? (\textquotedblleft{}A mother's brother\textquotedblright{}.) Say it once more.
\end{tcolorbox}
\section[\texorpdfstring{\pa{ਮਾਸੀ} (māsī)}{ਮਾਸੀ (māsī)}]{\pa{ਮਾਸੀ} (māsī) — a mother's sister}

\label{lesson:PA-C331-masi}



\begin{quote}
Before the new one: say the Punjabi for a father's elder brother, then the Punjabi for a mother's brother.
\end{quote}

\subsection*{You'll want to know: \pa{ਮਾਸੀ}}
\textbf{\pa{ਮਾਸੀ}} — \emph{māsī} — \textquotedblleft{}a mother's sister\textquotedblright{}.

\begin{quote}
\pa{ਮੇਰੀ} \pa{ਮਾਸੀ} \pa{ਹਰ} \pa{ਸਾਲ} \pa{ਸਾਡੇ} \pa{ਲਈ} \pa{ਤੋਹਫ਼ੇ} \pa{ਲਿਆਉਂਦੀ} \pa{ਹੈ।}
\end{quote}

\begin{quote}
\emph{merī māsī har sāl sāḍe laī tohfe liāundī hai.}
\end{quote}

\begin{quote}
My mother's sister brings us gifts every year.
\end{quote}

A feminine noun.

\begin{cousinweb}
\textbf{\pa{ਮਾਸੀ}} (\emph{māsī}) is inherited from Sanskrit \emph{mātuḥṣvasṛ}, \textquotedblleft{}a mother's sister\textquotedblright{}; CDIAL lists Punjabi \emph{māssī}, \textquotedblleft{}mother's sister\textquotedblright{}, under it (CDIAL 10001).
\end{cousinweb}

\begin{grammarlens}[title={What you've built}]
One more word for this chapter.
\end{grammarlens}

\subsection*{Your turn}
\begin{itemize}
\raggedright
  \item \emph{Say it:} \emph{māsī}
  \item \emph{Say it:} \emph{māsī}, once more
  \item \emph{Say it:} \emph{merī māsī har sāl sāḍe laī tohfe liāundī hai.} — My mother's sister brings us gifts every year.
  \item \emph{Recall:} say the Punjabi for a father's elder brother, then the Punjabi for a mother's brother, then say \emph{māsī} again
  \item \emph{Write it:} copy \textbf{\pa{ਮਾਸੀ}} once, checking the whole word
\end{itemize}

\begin{tcolorbox}[breakable,colback=teal!4,colframe=teal!35!black,title={Before you move on}]
What does \pa{ਮਾਸੀ} mean? (\textquotedblleft{}A mother's sister\textquotedblright{}.) Say it once more.
\end{tcolorbox}
\section[\texorpdfstring{\pa{ਮਾਸੜ} (māsaṛ)}{ਮਾਸੜ (māsaṛ)}]{\pa{ਮਾਸੜ} (māsaṛ) — a mother's sister's husband}

\label{lesson:PA-C331-masar}



\begin{quote}
Before the new one: say the Punjabi for a mother's brother, then the Punjabi for a mother's sister.
\end{quote}

\subsection*{You'll want to know: \pa{ਮਾਸੜ}}
\textbf{\pa{ਮਾਸੜ}} — \emph{māsaṛ} — \textquotedblleft{}a mother's sister's husband\textquotedblright{}.

\begin{quote}
\pa{ਮੇਰਾ} \pa{ਮਾਸੜ} \pa{ਇੱਕ} \pa{ਦਰਜ਼ੀ} \pa{ਹੈ।}
\end{quote}

\begin{quote}
\emph{merā māsaṛ ikk darzī hai.}
\end{quote}

\begin{quote}
My mother's sister's husband is a tailor.
\end{quote}

A masculine noun.

\begin{cousinweb}
\textbf{\pa{ਮਾਸੜ}} (\emph{māsaṛ}) is inherited from Sanskrit \emph{mātuḥṣvasṛ}, \textquotedblleft{}a mother's sister\textquotedblright{}; CDIAL lists Punjabi \emph{māsaṛ}, \textquotedblleft{}mother's sister's husband\textquotedblright{}, under it (CDIAL 10001).
\end{cousinweb}

\begin{grammarlens}[title={What you've built}]
One more word for this chapter.
\end{grammarlens}

\subsection*{Your turn}
\begin{itemize}
\raggedright
  \item \emph{Say it:} \emph{māsaṛ}
  \item \emph{Say it:} \emph{māsaṛ}, once more
  \item \emph{Say it:} \emph{merā māsaṛ ikk darzī hai.} — My mother's sister's husband is a tailor.
  \item \emph{Recall:} say the Punjabi for a mother's brother, then the Punjabi for a mother's sister, then say \emph{māsaṛ} again
  \item \emph{Write it:} copy \textbf{\pa{ਮਾਸੜ}} once, checking the whole word
\end{itemize}

\begin{tcolorbox}[breakable,colback=teal!4,colframe=teal!35!black,title={Before you move on}]
What does \pa{ਮਾਸੜ} mean? (\textquotedblleft{}A mother's sister's husband\textquotedblright{}.) Say it once more.
\end{tcolorbox}
